% Options for packages loaded elsewhere
\PassOptionsToPackage{unicode}{hyperref}
\PassOptionsToPackage{hyphens}{url}
\PassOptionsToPackage{dvipsnames,svgnames,x11names}{xcolor}
%
\documentclass[
  11pt,
]{article}
\usepackage{amsmath,amssymb}
\usepackage{iftex}
\ifPDFTeX
  \usepackage[T1]{fontenc}
  \usepackage[utf8]{inputenc}
  \usepackage{textcomp} % provide euro and other symbols
\else % if luatex or xetex
  \usepackage{unicode-math} % this also loads fontspec
  \defaultfontfeatures{Scale=MatchLowercase}
  \defaultfontfeatures[\rmfamily]{Ligatures=TeX,Scale=1}
\fi
\usepackage{lmodern}
\ifPDFTeX\else
  % xetex/luatex font selection
\fi
% Use upquote if available, for straight quotes in verbatim environments
\IfFileExists{upquote.sty}{\usepackage{upquote}}{}
\IfFileExists{microtype.sty}{% use microtype if available
  \usepackage[]{microtype}
  \UseMicrotypeSet[protrusion]{basicmath} % disable protrusion for tt fonts
}{}
\makeatletter
\@ifundefined{KOMAClassName}{% if non-KOMA class
  \IfFileExists{parskip.sty}{%
    \usepackage{parskip}
  }{% else
    \setlength{\parindent}{0pt}
    \setlength{\parskip}{6pt plus 2pt minus 1pt}}
}{% if KOMA class
  \KOMAoptions{parskip=half}}
\makeatother
\usepackage{xcolor}
\usepackage[left=1.25in,right=1.25in,top=1in,bottom=1in]{geometry}
\usepackage{longtable,booktabs,array}
\usepackage{calc} % for calculating minipage widths
% Correct order of tables after \paragraph or \subparagraph
\usepackage{etoolbox}
\makeatletter
\patchcmd\longtable{\par}{\if@noskipsec\mbox{}\fi\par}{}{}
\makeatother
% Allow footnotes in longtable head/foot
\IfFileExists{footnotehyper.sty}{\usepackage{footnotehyper}}{\usepackage{footnote}}
\makesavenoteenv{longtable}
\setlength{\emergencystretch}{3em} % prevent overfull lines
\providecommand{\tightlist}{%
  \setlength{\itemsep}{0pt}\setlength{\parskip}{0pt}}
\setcounter{secnumdepth}{-\maxdimen} % remove section numbering
\ifLuaTeX
  \usepackage{selnolig}  % disable illegal ligatures
\fi
\IfFileExists{bookmark.sty}{\usepackage{bookmark}}{\usepackage{hyperref}}
\IfFileExists{xurl.sty}{\usepackage{xurl}}{} % add URL line breaks if available
\urlstyle{same}
\hypersetup{
  pdftitle={GPS satellite relativistic corrections --- standard derivation, proper-time companion, and open questions for a curved-spacetime extension},
  pdfauthor={Trey Morris with Claude Opus 4.7},
  colorlinks=true,
  linkcolor={blue},
  filecolor={Maroon},
  citecolor={Blue},
  urlcolor={blue},
  pdfcreator={LaTeX via pandoc}}

\title{GPS satellite relativistic corrections --- standard derivation,
proper-time companion, and open questions for a curved-spacetime
extension}
\author{Trey Morris with Claude Opus 4.7}
\date{2026-06-01}

\begin{document}
\maketitle

\hypertarget{cover-note}{%
\section{§1 Cover note}\label{cover-note}}

This is a short follow-up to the interim report of
2026-05-25\footnote{Interim report of 2026-05-25.
  \texttt{Roadmapping/Author\_Reports/2026-05\_interim\_for\_gill.md}.}.
We narrow it to one applied-relativity question that opened in the
public repository the same week, namely the relativistic
clock-correction problem for GPS satellites. The question was prompted
by a popular-science article on the \(\pm 38\,\mu\)s/day
offset\footnote{Issue \#57 --- GPS relativistic clock-correction
  question. \texttt{github.com/temoTxt/PyPhysics/issues/57}.}. Two
campaigns followed. The first\footnote{Standard-method derivation
  campaign. PR \#71, \texttt{github.com/temoTxt/PyPhysics/pull/71}.} is
a standard-method derivation reproducing Ashby's \emph{Living Reviews}
treatment\footnote{Ashby, N. (2003). \emph{Relativity in the Global
  Positioning System.} Living Reviews in Relativity.} end-to-end. The
second\footnote{Proper-time companion campaign. PR \#73,
  \texttt{github.com/temoTxt/PyPhysics/pull/73}.} is a proper-time
companion that applies the Gill--Zachary substitution rules\footnote{Gill,
  T. L. and Zachary, W. W. \emph{Two Mathematically Equivalent Versions
  of Maxwell's Equations.} See
  \texttt{Roadmapping/Equation\_Verification/Two\_Mathematically\_Equivalent\_Versions\_of\_Maxwells\_Equations.md}
  for the per-equation verification of the substitution rules used here.}
to each of the nine effects. The companion campaign surfaced four open
questions about how the framework should extend to curved spacetime. The
framework as-published does not yet answer them. This report
consolidates the two campaigns and the four questions in a form short
enough to read alongside the broader interim report.

\hypertarget{what-this-report-is-and-is-not}{%
\subsection{What this report is and is
not}\label{what-this-report-is-and-is-not}}

We intend this as a \emph{consolidation} of two campaigns whose
underlying work is committed to the public repository in full. It is not
a re-derivation. It is not a journal-style summary. It is not a request
that you read the per-effect documents to interpret the questions in §5.
Each question is self-contained. The load-bearing section is §5.
Everything else is context to make §5 readable in one sitting.

\hypertarget{the-gps-question}{%
\section{§2 The GPS question}\label{the-gps-question}}

GPS satellites carry atomic clocks. Once in orbit, those clocks tick
faster than ground clocks by

\[+38.57\,\mu\text{s/day}, \qquad (1)\]

a combined gain decomposing into a gravitational-redshift contribution
of \(+45.65\,\mu\)s/day and a special-relativistic velocity-dilation
contribution of \(-7.11\,\mu\)s/day. The engineering response is to slow
the satellite clock by exactly this fractional rate before launch. We
observe that the fundamental clock frequency of \(10.230\,000\,000\,00\)
MHz is built at \(10.229\,999\,995\,43\) MHz, a fractional offset

\[\Delta f / f = -4.4647 \times 10^{-10}. \qquad (2)\]

Once in orbit, the relativistic gain exactly cancels the pre-launch
slowdown. The satellite clock then matches a ground clock at the geoid
potential. This is the most stringently-tested applied-relativity result
available. Indeed, the satellite constellation has operated since 1977
with the relativistic correction validated daily against the
\(\sim 10^{-10}\) precision of cesium and rubidium clocks. Ashby (2003)
is the canonical reference\footnote{Ashby, N. (2003). \emph{Relativity
  in the Global Positioning System.} Living Reviews in Relativity.}; the
IAU 2000 geoid potential supplies the ground-side reference, and the
operational pre-launch offset is recorded in the GPS Interface Control
Document\footnote{GPS Interface Control Document ICD-200. Operational
  pre-launch frequency offset; IERS / WGS-84 / EGM2008 supply geoid and
  Earth-model parameters.}.

The proper-time framework's perspective is operationally aligned with
how GPS actually works. The satellite's atomic clock \emph{is} its
proper time \(\tau\). The question ``how does \(\tau_{\rm sat}\) relate
to \(\tau_{\rm grnd}\)?'' is exactly the kinematic question the
framework's apparatus

\[b = \sqrt{c^2 + \mathbf{u}^2}, \qquad \frac{dt}{d\tau} = \frac{b}{c} \qquad (3)\]

is designed to answer. The campaign set out to determine, per effect,
whether the framework reproduces standard SR + GR at GPS-relevant
precision. The expected outcome, given the framework's claim of
mathematical equivalence\footnote{Gill, T. L. and Zachary, W. W.
  \emph{Two Mathematically Equivalent Versions of Maxwell's Equations.}
  See
  \texttt{Roadmapping/Equation\_Verification/Two\_Mathematically\_Equivalent\_Versions\_of\_Maxwells\_Equations.md}
  for the per-equation verification of the substitution rules used here.},
is that it does. The campaign also sought to surface any structural
difference that would be operationally distinguishable at the
next-generation optical-clock precision floor (of order \(10^{-17}\)),
or in regimes the framework as-published does not cover.

\hypertarget{standard-derivation-what-the-first-campaign-reproduces}{%
\section{§3 Standard derivation --- what the first campaign
reproduces}\label{standard-derivation-what-the-first-campaign-reproduces}}

The standard campaign\footnote{Standard-method derivation campaign. PR
  \#71, \texttt{github.com/temoTxt/PyPhysics/pull/71}.} covered nine
separately-named effects. Each is documented in its own per-effect file
with a Wolfram-MCP numerical check and a comparison against Ashby
(2003)\footnote{Ashby, N. (2003). \emph{Relativity in the Global
  Positioning System.} Living Reviews in Relativity.}. The combined
decomposition reproduces Ashby end-to-end at the precision the published
parameters can sustain.

\begin{longtable}[]{@{}
  >{\raggedright\arraybackslash}p{(\columnwidth - 4\tabcolsep) * \real{0.3333}}
  >{\raggedright\arraybackslash}p{(\columnwidth - 4\tabcolsep) * \real{0.3333}}
  >{\raggedright\arraybackslash}p{(\columnwidth - 4\tabcolsep) * \real{0.3333}}@{}}
\toprule\noalign{}
\begin{minipage}[b]{\linewidth}\raggedright
Effect
\end{minipage} & \begin{minipage}[b]{\linewidth}\raggedright
Standard-method magnitude
\end{minipage} & \begin{minipage}[b]{\linewidth}\raggedright
Ashby cross-reference
\end{minipage} \\
\midrule\noalign{}
\endhead
\bottomrule\noalign{}
\endlastfoot
02 Gravitational time dilation & \(+45.65\,\mu\)s/day & §3.1 Eq. (19) \\
03 Velocity time dilation & \(-7.11\,\mu\)s/day & §3.2 Eq. (20) \\
04 Combined offset including \(J_2\) & \(+38.57\,\mu\)s/day & §3.3 Eq.
(39) \\
Pre-launch frequency offset & \(\Delta f/f = -4.4647 \times 10^{-10}\) &
§3.3 Eq. (39) \\
05 Eccentricity periodic (\(e = 0.02\)) & \(\pm 46\) ns peak & §4 Eq.
(43) \\
06 Sagnac maximum & \(\sim 137\) ns per signal & §3.4 Eq. (35) \\
07 Shapiro maximum & \(\sim 62\) ps per signal & §6 Eq. (53) \\
08 \(J_2\) contribution to geoid & \(+0.033\,\mu\)s/day & §3.5 Eq.
(24) \\
\end{longtable}

The campaign's value to the repository is \emph{consolidation}. We
collect Ashby's per-effect derivations into one place with Wolfram-MCP
verification of every quoted number against named parameter
values\footnote{Parameter sources: IERS conventions, WGS-84, EGM2008,
  GPS ICD-200.}. This is a reproduction-of-standard exercise. No novel
claim is made. We flag the campaign here because §4's proper-time
companion is conditional on the §3 baseline being right.

\hypertarget{proper-time-companion-what-the-second-campaign-found}{%
\section{§4 Proper-time companion --- what the second campaign
found}\label{proper-time-companion-what-the-second-campaign-found}}

The companion campaign\footnote{Proper-time companion campaign. PR \#73,
  \texttt{github.com/temoTxt/PyPhysics/pull/73}.} re-derived the same
nine effects under the Gill--Zachary substitution rules\footnote{Gill,
  T. L. and Zachary, W. W. \emph{Two Mathematically Equivalent Versions
  of Maxwell's Equations.} See
  \texttt{Roadmapping/Equation\_Verification/Two\_Mathematically\_Equivalent\_Versions\_of\_Maxwells\_Equations.md}
  for the per-equation verification of the substitution rules used here.},

\[\frac{\mathbf{w}}{c} = \frac{\mathbf{u}}{b}, \qquad \frac{1}{c}\frac{\partial}{\partial t} = \frac{1}{b}\frac{\partial}{\partial \tau}, \qquad b^2 = c^2 + \mathbf{u}^2. \qquad (4)\]

We recorded, per effect, whether the framework applies as-published or
requires a speculative extension. The nine per-effect documents split
into three categories.

\textbf{Pass --- direct framework application (pt\_03, pt\_06).} Two
effects admit the verified substitution rules without modification. The
most direct demonstration is pt\_03 velocity time dilation. We observe
that the framework's identity

\[\frac{d\tau}{dt} = \frac{c}{b} = \sqrt{1 - \mathbf{w}^2/c^2} = \frac{1}{\gamma} \qquad (5)\]

follows by algebraic identity from the velocity-duality substitution of
(4). The numerical value \(-7.11\,\mu\)s/day is identical to the
standard prediction at all decimal places, not at any precision floor.
For pt\_06 Sagnac, the rotating-frame photon-null condition runs through
with \(c \to b\) substitution. The result is

\[\Delta t = -\frac{2\boldsymbol{\omega} \cdot \mathbf{A}}{b^2}, \qquad (6)\]

which reduces to the standard
\(-2\boldsymbol{\omega} \cdot \mathbf{A} / c^2\) at \(b \approx c\) for
an Earth-rotating receiver. The residual sits at \(\sim 10^{-10}\) ns,
invisible to current GPS clocks and marginal at next-generation
optical-clock precision.

\textbf{Marginal --- kinematic Pass plus speculative GR extension
(pt\_01, pt\_04, pt\_05).} Three effects admit the kinematic piece as
direct framework application but require a curved-spacetime extension
that the framework as-published does not provide. The campaign adopts a
\emph{minimal-extension hypothesis}: replace \(c\) by the local \(b\) in
the Schwarzschild line element. Under that hypothesis, the predictions
reduce to standard at GPS-relevant velocities. The combined offset of
(1) is reproduced. The operational pre-launch offset of (2) is
unchanged. The pt\_05 eccentricity case additionally surfaces a
framework-specific third-term contribution,

\[\frac{\mathbf{u} \cdot \mathbf{a}}{b^4} \, \tau_{\rm signal}, \qquad (7)\]

arising from the wave-equation dissipative coefficient in Maxwell-paper
Eq. (4)\footnote{Gill, T. L. and Zachary, W. W. \emph{Two Mathematically
  Equivalent Versions of Maxwell's Equations.} See
  \texttt{Roadmapping/Equation\_Verification/Two\_Mathematically\_Equivalent\_Versions\_of\_Maxwells\_Equations.md}
  for the per-equation verification of the substitution rules used here.}.
For GPS circular orbits this vanishes (\(\mathbf{u} \perp \mathbf{a}\)).
For elliptical orbits at \(e = 0.02\) it is of order \(10^{-34}\) s per
signal, operationally invisible. We note that the same coefficient
appears in the radiation-reaction sub-investigation\footnote{Radiation-reaction
  sub-investigation. Issue \#43,
  \texttt{github.com/temoTxt/PyPhysics/issues/43}.}, where it
contributes at order unity. The framework is thus internally consistent
on which structures appear where.

\textbf{Out of scope --- framework as-published does not extend (pt\_02,
pt\_07, pt\_08).} Three effects are purely gravitational: the
Schwarzschild redshift, the Shapiro propagation delay, and the \(J_2\)
multipole. The framework's substitution rules govern SR Maxwell and
flat-space dynamics; they do not extend to curved spacetime in any
verified paper. Under the same minimal \(c \to b\) Schwarzschild
extension, each effect reduces to standard at GPS precision. However,
the speculative extension is the load-bearing assumption, not a derived
result. We therefore record these three as out-of-scope for the
as-published framework, not as refuted by it. Every document in this
category carries a substantive-AI TODO block on its §0 ``Framework
applicability'' paragraph.

\textbf{Operational summary.} At GPS precision
(\(\mathbf{u}^2/c^2 \approx 1.7 \times 10^{-10}\); gravitational
\(GM/(rc^2) \approx 10^{-10}\); clock stability \(\sim 10^{-14}\)), the
framework reproduces the standard prediction for every effect. The
pre-launch frequency offset of (2) is unchanged; the operational
engineering response is unchanged; no observable deviation was found.
This is the expected outcome of mathematical equivalence in the
appropriate limit; it is consistent with, but does not independently
validate, either formulation. We observe that any disagreement with
measurement at GPS precision would have been a serious problem for the
framework; none was found.

\hypertarget{four-open-questions-for-a-curved-spacetime-extension}{%
\section{§5 Four open questions for a curved-spacetime
extension}\label{four-open-questions-for-a-curved-spacetime-extension}}

This is the load-bearing section of the report. The four questions all
surface from the proper-time companion's Marginal and Out-of-scope
per-effect documents. Each is operationally invisible at
current-generation GPS precision. Each is structurally load-bearing for
the framework's extension to regimes where the differences would matter.
Each question is one short paragraph by design. The per-effect documents
under the GPS-relativity folder\footnote{Per-effect documents.
  \texttt{Roadmapping/GPS\_Relativity/}.} carry the supporting context.

\textbf{Q1. What is the correct curved-spacetime extension of the
framework?} The minimal-extension hypothesis used in pt\_01, pt\_02,
pt\_04, pt\_07, and pt\_08 is \(c \to b\) everywhere in the
Schwarzschild line element, with \(b \to c\) at low velocity. This is
one possible extension. We considered three others without resolution:

\begin{itemize}
\item
  \begin{enumerate}
  \def\labelenumi{(\alph{enumi})}
  \tightlist
  \item
    the substitution \(c \to b\) applied only in the time-time metric
    component \(g_{00}\);
  \end{enumerate}
\item
  \begin{enumerate}
  \def\labelenumi{(\alph{enumi})}
  \setcounter{enumi}{1}
  \tightlist
  \item
    treating null and timelike geodesics differently, since the photon's
    local frame has \(\mathbf{u} = 0\) and so \(b = c\), but the
    worldline along which the metric is evaluated does not;
  \end{enumerate}
\item
  \begin{enumerate}
  \def\labelenumi{(\alph{enumi})}
  \setcounter{enumi}{2}
  \tightlist
  \item
    replacing the metric structure entirely with a \(b\)-dependent
    connection that reduces to the Schwarzschild connection in the
    \(b = c\) limit.
  \end{enumerate}
\end{itemize}

For GPS at the \(10^{-10}\) precision floor, all four hypotheses reduce
to the same operational answer. The question matters for next-generation
optical-clock GPS, where \(10^{-17}\) stability is at the edge of
distinguishing them, and for strong-field regimes where \(GM/(rc^2)\) is
no longer of order \(10^{-10}\). The campaign cannot resolve this from
the verified corpus. We flag it as the principal Tepper-input ask of
this report.

\textbf{Q2. How does \(\mathbf{u} \cdot \mathbf{a}\) transform under the
proper-time boost?} The boost in question is the one from Maxwell-paper
Eq. (11)\footnote{Gill, T. L. and Zachary, W. W. \emph{Two
  Mathematically Equivalent Versions of Maxwell's Equations.} See
  \texttt{Roadmapping/Equation\_Verification/Two\_Mathematically\_Equivalent\_Versions\_of\_Maxwells\_Equations.md}
  for the per-equation verification of the substitution rules used here.}.
This is the same open question recorded in the proper-time cheat-sheet's
open-questions list\footnote{Proper-time cheat-sheet, open-questions
  list §4a.
  \texttt{Roadmapping/Electromagnetism/\_proper\_time\_cheatsheet.md}.},
flagged for the radiation-reaction sub-investigation\footnote{Radiation-reaction
  sub-investigation. Issue \#43,
  \texttt{github.com/temoTxt/PyPhysics/issues/43}.}. It surfaces again
in pt\_05 of the GPS campaign. We observe that the third-term
contribution of (7) depends on \(\mathbf{u} \cdot \mathbf{a}\) having a
frame-independent operational meaning. The covariance of this 3-vector
dot product under the proper-time boost has not been derived in any
verified paper. For GPS the question is moot
(\(\mathbf{u} \cdot \mathbf{a} = 0\) for circular orbits; of order
\(10^{-34}\) s for elliptical). However, it bears on whether the
radiation-reaction predictions of the sub-investigation\footnote{Radiation-reaction
  sub-investigation. Issue \#43,
  \texttt{github.com/temoTxt/PyPhysics/issues/43}.} are
frame-independent at the level the experiments need. Even a one-line
confirmation that \(\mathbf{u} \cdot \mathbf{a}\) transforms as a
Lorentz scalar (or, conversely, that it does not) would resolve this for
both campaigns.

\textbf{Q3. Does the framework predict a satellite-side Sagnac
contribution?} The standard Sagnac correction
\(\Delta t = -2\boldsymbol{\omega} \cdot \mathbf{A} / c^2\) carries
\(c\) because it is derived in the rotating frame of the
\emph{receiver}. Under the proper-time framework, the receiver's
effective light-speed is

\[b_{\rm recv} = \sqrt{c^2 + u_{\rm recv}^2}, \qquad (8)\]

where \(u_{\rm recv} \leq 465\) m/s is the ground rotation, giving
\(b_{\rm recv}^2/c^2 - 1 \approx 1.2 \times 10^{-12}\). The
receiver-side derivation in pt\_06 uses \(b_{\rm recv}\) and recovers
the standard prediction at this precision. A symmetric reading of the
framework would have the satellite's

\[b_{\rm sat} = \sqrt{c^2 + u_{\rm sat}^2} \qquad (9)\]

enter the Sagnac formula on the source side as well, with
\(u_{\rm sat} \approx 3874\,\text{m/s}\) giving
\(b_{\rm sat}^2/c^2 - 1 \approx 1.7 \times 10^{-10}\). That reading
would predict an additional satellite-side correction at the
\(\sim 10^{-10}\) level of the standard Sagnac. This sits below current
GPS precision, but at the edge of next-generation optical-clock
receivers. We ask whether the framework intends the source-side and
receiver-side \(b\) to enter symmetrically, or whether the
receiver-frame derivation is the complete prediction.

\textbf{Q4. Are there framework-specific \(J_2\)-like corrections beyond
the speculative extension?} Earth's quadrupole moment
\(J_2 = 1.083 \times 10^{-3}\) modifies the gravitational potential
through a \(P_2(\cos\theta)\) multipole term. Under the speculative
\(c \to b\) extension, the \(J_2\) correction enters via the same
\(1/b^2\) factor as the leading point-mass term. It reproduces the
standard \(+0.033\,\mu\)s/day contribution to the combined offset. A
first-principles derivation of multipole-induced clock corrections from
the framework's underlying dynamics would clarify whether the standard
contribution is the leading-order term, or whether additional
\(b\)-dependent contributions arise at higher multipole order. The
Lense--Thirring (frame-dragging) contribution from Earth's rotation sits
in the same category. It is at the \(\sim 10^{-16}\) level for current
GPS, but within reach of ACES on the ISS. The framework's prediction
here would need a curved-spacetime gravitomagnetic extension that we
have not derived. For GPS the Q4 effect is sub-picosecond per day. For
next-generation Earth-orbiting clock missions it is the leading
post-Schwarzschild correction.

\hypertarget{suggested-next-steps}{%
\section{§6 Suggested next steps}\label{suggested-next-steps}}

We order these by what would unblock the most downstream work.

\begin{enumerate}
\def\labelenumi{\arabic{enumi}.}
\item
  \textbf{Resolve Q1.} Of the four questions, Q1 is the one that
  determines whether the others have well-defined answers; Q3 and Q4 are
  partially dependent on it, and Q2 is independent and propagates to the
  radiation-reaction sub-investigation\footnote{Radiation-reaction
    sub-investigation. Issue \#43,
    \texttt{github.com/temoTxt/PyPhysics/issues/43}.} regardless. Even a
  directional answer would suffice. Such an answer might take the form
  ``the right extension is \(c \to b\) everywhere'', or ``it is the
  \(c \to b\) only in \(g_{00}\) variant'', or ``the framework's
  intended GR extension is not yet written down, but the right answer is
  the cleanest minimal coupling consistent with the Maxwell paper's
  substitution rules''. Any of these would let the GPS campaign convert
  the affected Marginal and Out-of-scope per-effect documents from
  speculative to derived.
\item
  \textbf{Address Q2 across two campaigns.} A
  \(\mathbf{u} \cdot \mathbf{a}\) transformation rule would close out
  the third-term contribution of (7) in both pt\_05 (GPS) and the
  radiation-reaction sub-investigation\footnote{Radiation-reaction
    sub-investigation. Issue \#43,
    \texttt{github.com/temoTxt/PyPhysics/issues/43}.}. The
  radiation-reaction sub-investigation is the regime where
  \(\mathbf{u} \cdot \mathbf{a}\) is operationally observable; the GPS
  context is incidental. The cross-campaign consistency is the
  load-bearing reason to write the rule down.
\item
  \textbf{Defer Q3 and Q4 to next-generation precision.} Neither
  question affects current GPS operations. We flag them here for
  completeness and for future-mission scope, namely optical-clock GPS,
  ACES, and mission concepts at the \(10^{-17}\)--\(10^{-18}\)
  clock-stability floor. A directional answer on either would be
  sufficient to scope follow-on threads.
\end{enumerate}

A natural next campaign, should the disposition of Q1 warrant it, is a
derived curved-spacetime extension of the framework --- pursued either
through the dual-Dirac equation or the proper-time photon propagator. We
would scope such a campaign as a follow-on issue, keyed off the
direction given. If, alternatively, the as-published framework's silence
on GR is intentional --- the framework is SR-and-flat-space by design,
leaving GR to the standard GR apparatus --- confirming that would let us
close the four open questions in this report with that disposition
recorded.

The prior interim report's §6 \(r_e\) question\footnote{Interim report
  of 2026-05-25.
  \texttt{Roadmapping/Author\_Reports/2026-05\_interim\_for\_gill.md}.}
remains the highest-load ask across all campaigns. This report's Q1 is
the highest-load ask within the GPS scope; it is, however, a much
smaller decision, since GPS is operationally non-falsifying at current
precision.

\hypertarget{references}{%
\subsection{References}\label{references}}

\end{document}
